\documentclass[10pt,a4paper]{article}

% ============================================================
% B2 - PROTOCOLO DE PROTECCION DE DATOS PERSONALES
% Documento autocontenido y compatible con pdfLaTeX en Overleaf.
% ============================================================

\usepackage[utf8]{inputenc}
\usepackage[T1]{fontenc}
\usepackage[spanish,es-nodecimaldot]{babel}
\usepackage{tgheros}
\renewcommand{\familydefault}{\sfdefault}
\usepackage[a4paper,top=1.85cm,bottom=1.8cm,left=1.9cm,right=1.9cm,headheight=20pt]{geometry}
\usepackage{xcolor}
\usepackage{tabularx,array,booktabs,colortbl}
\usepackage{enumitem}
\usepackage[most]{tcolorbox}
\usepackage{fancyhdr}
\usepackage{lastpage}
\usepackage{microtype}
\usepackage{setspace}
\usepackage{titlesec}
\usepackage{hyperref}
\usepackage{xurl}
\usepackage{tikz}
\usepackage{ragged2e}

% ---------- PALETA ROUTTECH ----------
\definecolor{RTForest}{HTML}{25463B}
\definecolor{RTTerracotta}{HTML}{B45F45}
\definecolor{RTGold}{HTML}{D5A64A}
\definecolor{RTCream}{HTML}{F7F2E8}
\definecolor{RTLight}{HTML}{EEF3F0}
\definecolor{RTGray}{HTML}{4D5652}

% ---------- DATOS ----------
\newcommand{\Proyecto}{RoutTech: Sistema de Movilidad Colaborativa para la Comunidad Universitaria de la UTEQ}
\newcommand{\Asignatura}{ISR-401 -- Ingeniería de Requisitos}
\newcommand{\Periodo}{2026--2027 PPA}
\newcommand{\CategoriaEtica}{Categoría B -- Datos personales}
\newcommand{\Repositorio}{https://github.com/AnderJM3/RoutTech_ISR401}
\newcommand{\FechaDocumento}{29 de julio de 2026}
\newcommand{\VersionDocumento}{2.0}
\newcommand{\Paralelo}{Cuarto nivel -- Software B}
\newcommand{\EstadoDocumento}{Versión 2.0 para revisión ética}
\newcommand{\Anderson}{NARANJO FLORES ANDERSON JEAMPIERE}
\newcommand{\Mayra}{CORDOVA CARREÑO MAYRA LUCILA}

% ---------- TABLAS ----------
\newcolumntype{Y}{>{\RaggedRight\arraybackslash}X}
\newcolumntype{C}[1]{>{\centering\arraybackslash}p{#1}}
\newcolumntype{L}[1]{>{\RaggedRight\arraybackslash}p{#1}}
\renewcommand{\arraystretch}{1.20}
\setlength{\tabcolsep}{4.5pt}

% ---------- ENCABEZADO Y PIE ----------
\pagestyle{fancy}
\fancyhf{}
\fancyhead[L]{\small\textbf{UTEQ | FCC | ISR-401}}
\fancyhead[R]{\small\textbf{B2 -- Protección de datos personales}}
\fancyfoot[L]{\small Categoría B -- RoutTech}
\fancyfoot[R]{\small Página \thepage\ de \pageref{LastPage}}
\renewcommand{\headrulewidth}{0.8pt}
\renewcommand{\footrulewidth}{0.4pt}
\renewcommand{\headrule}{\hbox to\headwidth{\color{RTTerracotta}\leaders\hrule height \headrulewidth\hfill}}
\renewcommand{\footrule}{\hbox to\headwidth{\color{RTGold}\leaders\hrule height \footrulewidth\hfill}}

% ---------- TITULOS ----------
\titleformat{\section}
  {\large\bfseries\color{RTForest}}
  {\thesection.}{0.45em}{}
  [\vspace{1pt}\color{RTGold}\titlerule]
\titlespacing*{\section}{0pt}{0.75em}{0.35em}

% ---------- CAJA ----------
\newtcolorbox{InfoBox}[1][]{
  enhanced,breakable,
  colback=RTLight,colframe=RTForest,
  boxrule=0.8pt,arc=2mm,
  left=2.6mm,right=2.6mm,top=1.6mm,bottom=1.6mm,
  title=#1,fonttitle=\bfseries,coltitle=white,
  colbacktitle=RTForest
}

\hypersetup{
  colorlinks=true,
  linkcolor=RTForest,
  urlcolor=RTTerracotta,
  pdftitle={B2 Protocolo de Proteccion de Datos Personales RoutTech},
  pdfauthor={Equipo RoutTech}
}
\urlstyle{same}
\setlength{\parindent}{1.1em}
\setlength{\parskip}{2.5pt}
\setstretch{1.10}
\setlist[itemize]{leftmargin=1.35em,itemsep=1pt,topsep=2pt}

\begin{document}
\justifying

% ======================= PORTADA =======================
\thispagestyle{empty}
\begin{tikzpicture}[remember picture,overlay]
  \fill[RTForest] (current page.north west) rectangle ([yshift=-4.6cm]current page.north east);
  \fill[RTTerracotta] ([yshift=-4.6cm]current page.north west) rectangle ([yshift=-4.95cm]current page.north east);
\end{tikzpicture}

\vspace*{0.65cm}
\begin{center}
{\color{white}\Large\bfseries UNIVERSIDAD TÉCNICA ESTATAL DE QUEVEDO}

\vspace{1.75cm}
{\large\bfseries\color{RTForest} Facultad de Ciencias de la Computación\\[4pt]}
{\normalsize\color{RTGray} Carrera de Ingeniería de Software}

\vspace{0.8cm}

{\Huge\bfseries\color{RTForest} PROTOCOLO DE PROTECCIÓN\\[5pt] DE DATOS PERSONALES\\[8pt]}
{\Large\bfseries\color{RTTerracotta} B2 -- Categoría B}

\vspace{1.0cm}

\begin{tcolorbox}[width=0.90\textwidth,colback=RTCream,colframe=RTGold,boxrule=1pt,arc=2mm]
\centering
{\large\bfseries \Proyecto}\\[7pt]
\textbf{Clasificación ética:} \CategoriaEtica\\
\textbf{Asignatura:} \Asignatura\\
\textbf{Periodo académico:} \Periodo
\end{tcolorbox}

\vfill
\begin{tabular}{rl}
\textbf{Documento:} & B2 -- Protocolo de protección de datos personales\\
\textbf{Versión:} & \VersionDocumento\\
\textbf{Estado:} & \EstadoDocumento\\
\textbf{Paralelo:} & \Paralelo\\
\textbf{Repositorio:} & \href{\Repositorio}{RoutTech\_ISR401}\\
\textbf{Fecha:} & \FechaDocumento
\end{tabular}

\vspace{0.75cm}
{\small Quevedo -- Los Ríos -- Ecuador}
\end{center}
\clearpage

% ======================= CONTENIDO =======================
\section{Alcance}

Este protocolo establece las medidas aplicables al tratamiento de datos personales obtenidos durante entrevistas, cuestionarios, observaciones y sesiones de validación del proyecto RoutTech. Su aplicación corresponde a \Anderson{} y \Mayra{}, únicos integrantes del equipo investigador con acceso autorizado a las evidencias del estudio.

\section{Minimización de datos}

\begin{itemize}
  \item Se recopilará únicamente la información necesaria para responder las preguntas de investigación y validar los requisitos del sistema.
  \item No se solicitarán contraseñas, placas vehiculares reales, domicilios exactos, ingresos, credenciales institucionales, historiales individuales de viaje ni ubicación en tiempo real.
  \item El cuestionario se analizará de manera anónima y las entrevistas se identificarán mediante códigos.
  \item Los participantes serán exclusivamente personas mayores de edad que acepten voluntariamente el consentimiento informado.
\end{itemize}

\section{Seudonimización}

Los nombres y demás identificadores directos serán sustituidos por códigos antes del análisis. La relación entre código e identidad, cuando exista, se almacenará separadamente y no se incorporará al repositorio del proyecto.

\rowcolors{2}{RTLight}{white}
\noindent\begin{tabularx}{\textwidth}{L{5.2cm} Y}
\rowcolor{RTForest}
\textcolor{white}{\textbf{Perfil}} & \textcolor{white}{\textbf{Formato de código}}\\
Estudiante universitario & EST-001, EST-002, EST-003, etc.\\
Conductor potencial & CON-001, CON-002, CON-003, etc.\\
Personal docente o administrativo & ADM-001, ADM-002, ADM-003, etc.\\
Usuario técnico & TEC-001, TEC-002, TEC-003, etc.\\
Usuario no técnico & NTEC-001, NTEC-002, NTEC-003, etc.\\
Evidencia de entrevista & EV-ENT-001, EV-ENT-002, EV-ENT-003, etc.\\
Cuestionario & EV-CUE-01.\\
\end{tabularx}

\section{Cifrado, almacenamiento y transferencia}

\begin{itemize}
  \item Los datos crudos, consentimientos, grabaciones y transcripciones identificables se conservarán en carpetas locales protegidas mediante cifrado AES-256 o una alternativa de seguridad equivalente.
  \item Las transferencias autorizadas se realizarán únicamente mediante conexiones protegidas con TLS 1.3 o una versión segura vigente.
  \item Las contraseñas de archivos cifrados no se enviarán junto con los archivos y se comunicarán por un canal separado.
  \item El repositorio \url{\Repositorio} se utilizará solo para documentos académicos, código, matrices y evidencias completamente anonimizadas. No se subirán datos crudos ni archivos con identificadores personales.
\end{itemize}

\section{Control de acceso}

El acceso a datos identificables estará limitado a \Anderson{} y \Mayra{}. Cada integrante utilizará credenciales personales y no compartirá copias con terceros. Si una persona deja de formar parte del equipo, su acceso deberá revocarse inmediatamente y las copias bajo su custodia deberán devolverse o eliminarse de forma segura.

\section{Retención y destrucción}

Los datos identificables se conservarán durante un máximo de 24 meses contados desde el cierre académico del proyecto. Cumplido ese plazo, se eliminarán las grabaciones, consentimientos digitalizados, tablas de correspondencia y demás archivos identificables mediante borrado seguro. Los resultados agregados, anonimizados o sintéticos podrán conservarse con fines académicos y de reproducibilidad.

\section{Derechos de los participantes}

Los participantes podrán solicitar acceso, rectificación, actualización, eliminación u oposición respecto de sus datos identificables, así como retirar su autorización mientras la información aún pueda asociarse con su identidad. Las solicitudes se recibirán a través de los correos institucionales de los investigadores y se atenderán conforme al consentimiento informado y al Plan de Gestión de Datos.

\section{Publicación y difusión}

\begin{itemize}
  \item No se publicarán nombres, cédulas, matrículas, correos, teléfonos, rostros, voces, placas, domicilios ni rutas individuales.
  \item Las ubicaciones se presentarán únicamente como campus, sector, barrio, parroquia o zona general.
  \item Las transcripciones, capturas y tablas serán revisadas antes de su inclusión en GitHub, informes, exposiciones o manuscritos.
  \item Los resultados se difundirán exclusivamente de forma agregada, seudonimizada o completamente disociada.
\end{itemize}

\begin{InfoBox}[Responsables del tratamiento]
\textbf{NARANJO FLORES ANDERSON JEAMPIERE}: coordinación, anonimización, integración documental y publicación controlada.\\
\textbf{CORDOVA CARREÑO MAYRA LUCILA}: recolección, codificación temática, validación y custodia inicial de evidencias.
\end{InfoBox}

\end{document}
